% =============================================================================
% Bloque 1: Apertura (slides 1–2)
% =============================================================================

% --- Slide 1: Portada ---
\maketitle

% --- Slide 2: Roadmap (lista con descripciones, sin tiempos) ---
\begin{frame}{Hoja de ruta}
\small
\setbeamercolor{enumerate item}{fg=uaiorange}
\begin{enumerate}\setlength\itemsep{4pt}
  \item \textbf{Introducción} — el régimen \emph{few-shot} y un escenario representativo donde este trabajo ayuda.

  \item \textbf{Problema} — escasez de datos, el marco clásico del re-muestreo y la pregunta central.

  \item \textbf{Estado del arte} — tres familias de aumentación, tres deficiencias persistentes y una década de evolución.

  \item \textbf{Hipótesis y objetivos} — qué afirmamos y cómo se va a verificar.

  \item \textbf{Metodología} — cinco filtros geométricos, ponderación suave y un pipeline desacoplado del LLM.

  \item \textbf{Resultados} — \alert{8.000+ configuraciones}, +2.25 pp sobre SMOTE y un hallazgo contra-intuitivo.

  \item \textbf{Discusión} — utilidad práctica, limitaciones honestas y trabajo futuro.
\end{enumerate}
\end{frame}
